\section{Algorithm Details}
\label{app:algo-details}

\subsection{Detailed Explanation of Algorithm~\ref{alg:full-information}}
\label{app:algo-1-details}

This section follows Algorithm~\ref{alg:full-information} and its proof in \cite[Section~3]{12WSNE}.

\subsubsection{Preprocessing}

\begin{enumerate}
  \item Compute a Nash equilibrium \((\mathbf{x}^*, \mathbf{y}^*)\) of the zero-sum game \((\mathcal{R}, -\mathcal{R})\); this can be solved in polynomial time via linear programming.

        By the properties of zero-sum games, \(\mathbf{x}^*\) is the row player's \textbf{maximin strategy}:
        \[
          \mathcal{R}(\mathbf{x}^*, \mathbf{y}^*) = \max_{\mathbf{x} \in \Delta^n} \min_{\mathbf{y} \in \Delta^n} \mathcal{R}(\mathbf{x}, \mathbf{y})
        \]
        Hence for every column strategy \(\mathbf{y}\), \(\mathcal{R}(\mathbf{x}^*, \mathbf{y}^*) \le \mathcal{R}(\mathbf{x}^*, \mathbf{y})\). \textbf{This property will be used several more times later.}

  \item Compute a Nash equilibrium \((\hat{\mathbf{x}}, \hat{\mathbf{y}})\) of the zero-sum game \((-\mathcal{C}, \mathcal{C})\). Similarly, for every row strategy \(\mathbf{x}\), \(\mathcal{C}(\hat{\mathbf{x}}, \hat{\mathbf{y}}) \le \mathcal{C}(\mathbf{x}, \hat{\mathbf{y}})\).
\end{enumerate}

We then case-analyze the relationship between \(\mathcal{R}(\mathbf{x}^*, \mathbf{y}^*)\) and \(\mathcal{C}(\hat{\mathbf{x}}, \hat{\mathbf{y}})\). By symmetry, we only analyze the case \textbf{\(\mathcal{R}(\mathbf{x}^*, \mathbf{y}^*) \ge \mathcal{C}(\hat{\mathbf{x}}, \hat{\mathbf{y}})\)}:

\subsubsection{The \texorpdfstring{\(\mathcal{R}(\mathbf{x}^*, \mathbf{y}^*) \ge \mathcal{C}(\hat{\mathbf{x}}, \hat{\mathbf{y}})\)}{R(x*,y*) >= C(xh,yh)} branch}

\paragraph{Case 1(a): \(\mathcal{R}(\mathbf{x}^*, \mathbf{y}^*) \le 1/2\)}
In this case, we simply output \((\hat{\mathbf{x}}, \mathbf{y}^*)\). Its correctness proof is given in Appendix~\ref{app:case-one-low}.

\paragraph{Case 1(b): ``one high, one low''}
In this case, \(\mathcal{R}(\mathbf{x}^*, \mathbf{y}^*) > 1/2\), and the following linear feasibility system has a solution \(\mathbf{x}'\):
\[
  \begin{cases}
    \sum_{i \in \supp(\mathbf{x}^*)} \mathbf{x}'(i) \cdot \mathcal{C}_{ij} \le \frac12, & \forall j \in [n]                 \\[4pt]
    \mathbf{x}'(i) \ge 0,                                                               & \forall i \in \supp(\mathbf{x}^*) \\[4pt]
    \sum_{i \in \supp(\mathbf{x}^*)} \mathbf{x}'(i) = 1
  \end{cases}
\]
This linear program can be solved, or declared infeasible, in polynomial time.
If it is feasible, we output \((\mathbf{x}', \mathbf{y}^*)\).

Its correctness proof is given in Appendix~\ref{app:case-one-low-high}.

\paragraph{Case 1(c): ``both high''}
In this case, neither of the previous two conditions holds; that is, we cannot bound a single player's best response by \(\le 1/2\).

Define the \textbf{restricted game} \((\mathcal{R}_{\mathbf{x}^*}, \mathcal{C}_{\mathbf{x}^*})\): the row player may only use actions in \(\supp(\mathbf{x}^*)\), while the column player may use all \(n\) actions.

One can show that every Nash equilibrium \((\mathbf{x}, \mathbf{y})\) of \((\mathcal{R}_{\mathbf{x}^*}, \mathcal{C}_{\mathbf{x}^*})\) satisfies:
\[
  \mathcal{R}(\mathbf{x}, \mathbf{y}) > \frac12,\;\; \mathcal{C}(\mathbf{x}, \mathbf{y}) > \frac12
\]
The contradiction argument is given in Appendix~\ref{app:case-one-both-high}.

Hence \((\mathbf{x}, \mathbf{y})\) is a \(1/2\)-WSNE of the original game. So we only need to compute such an \((\mathbf{x}, \mathbf{y})\). However, we cannot in general do this in polynomial time using known methods, because computing a NE of a bimatrix game is \PPAD-complete \cite{DGP09}.

Here we use a different sampling theorem that preserves the equilibrium payoffs on the sampled supports:

\begin{lemma}[{Constant-Support Well-Supporting Lemma \cite[Theorem~2]{cfj14}}]
  For any Nash equilibrium \((\mathbf{x}, \mathbf{y})\) of the restricted game \((\mathcal{R}_{\mathbf{x}^*}, \mathcal{C}_{\mathbf{x}^*})\), there exists a \(\kappa(\delta)\)-uniform strategy profile \((\mathbf{w}, \mathbf{z})\) such that:
  \[
    \forall i \in \supp(\mathbf{w}),\; \mathcal{R}(e_i, \mathbf{z}) \ge \mathcal{R}(\mathbf{x}, \mathbf{y}) - \delta > \frac12 - \delta
  \]
  \[
    \forall j \in \supp(\mathbf{z}),\; \mathcal{C}(\mathbf{w}, e_j) \ge \mathcal{C}(\mathbf{x}, \mathbf{y}) - \delta > \frac12 - \delta
  \]
  Since the maximum possible payoff is \(1\) (the payoff matrices are normalized), the above conditions imply:
  \[
    \forall i \in \supp(\mathbf{w}),\ \mathcal{R}(e_i, \mathbf{z}) \ge \left(\frac12 - \delta\right) - \left(1 - \max_k \mathcal{R}(e_k, \mathbf{z})\right) = \max_k \mathcal{R}(e_k, \mathbf{z}) - \left(\frac12 + \delta\right)
  \]
  and similarly for the column player. Therefore, \((\mathbf{w}, \mathbf{z})\) is a \((1/2+\delta)\)-WSNE.
\end{lemma}

Here \(\kappa(\delta) = \left\lceil \frac{2 \ln (1/\delta)}{\delta^2} \right\rceil\) is a constant for fixed \(\delta \in (0,1)\). A \(k\)-uniform strategy profile \((\mathbf{w}, \mathbf{z})\) satisfies that every component of \(\mathbf{w}, \mathbf{z}\) is a non-negative integer multiple of \(\frac{1}{k}\), so it naturally follows that \(\max(|\supp(\mathbf{w})|, |\supp(\mathbf{z})|) \le \kappa(\delta)\), i.e., both supports have size at most \(\kappa(\delta)\). There are \(O(n^{\kappa(\delta)})\) such strategies per player and hence \(O(n^{2\kappa(\delta)})\) profiles. Since the lemma guarantees existence, it suffices to enumerate all profiles and check whether each one is a \((1/2+\delta)\)-WSNE. When \(\delta\) is a constant, the running time of the algorithm is polynomial in \(n\).

\begin{remark}
  In Lemma~\ref{lem:lmm_thm}, the sampling size \(k = O\left(\frac{\log n}{\varepsilon^2}\right)\) depends on \(n\), while here \(\kappa(\delta)\) depends only on constant \(\delta\). This is (probably) because in this case we only need to constrain the sampled support, whereas the LMM lemma needs to constrain all \(n\) pure strategies.
\end{remark}

\subsection{Detailed Explanation of Algorithm~\ref{alg:query-efficient}}
\label{app:algo-2-details}

This section follows the payoff-query variant in \cite[Section~4]{12WSNE}, with the probability-accounting correction noted above.

\subsubsection{Motivation}
In practice, the cost of providing the entire payoff matrix may be unacceptable; only the payoff of a pure strategy profile \((e_i, e_j)\) can be queried. In this setting, the goal of the algorithm is to output an approximate WSNE with as few queries as possible.

If we keep using Algorithm 1, we face two problems:
\begin{enumerate}
  \item The preprocessing requires exact Nash equilibria of zero-sum games, but without access to the full matrices \(\mathcal{R}, \mathcal{C}\) we can only compute approximate equilibria;
  \item Cases 1(b) and 1(c) require the full information of the restricted game \((\mathcal{R}_{\mathbf{x}^*}, \mathcal{C}_{\mathbf{x}^*})\), but the support of \(\mathbf{x}^*\) may be large, making the number of queries unacceptable.
\end{enumerate}

\subsubsection{Modifications}

\paragraph{Approximate equilibria of zero-sum games}

For \(n \times n\) zero-sum bimatrix games, there is a randomized algorithm that, with \(O\left(\frac{n \log n}{\varepsilon^2}\right)\) queries, outputs an \(\varepsilon\)-NE with success probability \(1 - n^{-1/8}\) \cite[Theorem~1]{goldbergRoth14}.

Fearnley and Savani extended this result to the \(\varepsilon\)-WSNE case:
\begin{lemma}[{Query-Efficient \(\varepsilon\)-WSNE in Zero-Sum Games \cite[Corollary~5.5]{fearnleySavani16}}]
  \label{lem:query-efficient-eps-wsne-in-zero-sum-games}
  With \(O\left(\frac{n \log n}{\varepsilon^4}\right)\) queries, one can output an \(\varepsilon\)-WSNE with success probability at least \(\left( 1 - n^{-\frac{1}{8}} \right)\left( 1 - \frac{2}{n} \right)^2\).
\end{lemma}

Now \((\mathbf{x}^*, \mathbf{y}^*)\) and \((\hat{\mathbf{x}}, \hat{\mathbf{y}})\) are no longer exact equilibria, and the relevant inequalities change accordingly.

By the definition of \(\varepsilon\)-WSNE:
\begin{align*}
  \mathcal{R}(\mathbf{x}^*, \mathbf{y}^*)         & \ge \mathcal{R}(e_i, \mathbf{y}^*) - \varepsilon     &  & \forall i \in [n] \\
  \mathcal{C}(\hat{\mathbf{x}}, \hat{\mathbf{y}}) & \ge \mathcal{C}(\hat{\mathbf{x}}, e_j) - \varepsilon &  & \forall j \in [n]
\end{align*}
Also, consider the corresponding form of the second inequality in the first game:
\begin{align*}
  -\mathcal{R}(\mathbf{x}^*, \mathbf{y}^*)                   & \ge -\mathcal{R}(\mathbf{x}^*, e_j) - \varepsilon &  & \forall j \in [n] \\
  \iff \mathcal{R}(\mathbf{x}^*, \mathbf{y}^*) - \varepsilon & \le \mathcal{R}(\mathbf{x}^*, e_j)                &  & \forall j \in [n]
\end{align*}
Similarly, we have:
\[
  \mathcal{C}(\hat{\mathbf{x}}, \hat{\mathbf{y}}) - \varepsilon \le \mathcal{C}(e_i, \hat{\mathbf{y}}) \;\; \forall i \in [n]
\]
Thus we obtain four inequalities for the approximate-NE setting. In the following, we call these four inequalities Observation 1 (they are Observation~2 in the original paper \cite[Observation~2]{12WSNE}).

\paragraph{Shrinking the support via the LMM lemma}

For the places in Algorithm 1 that need supports, taking \(\mathbf{x}^*\) as an example, we sample a strategy \(\mathbf{x}_s^*\) from \(\mathbf{x}^*\) to replace it. Since the support of \(\mathbf{x}_s^*\) has size \(O(\log n / \varepsilon^2)\), we can \textbf{query all} pairs \((e_i, e_j)\) with \(i \in \supp(\mathbf{x}_s^*)\) and \(j \in [n]\); the total number of such queries is \(O(n \cdot \frac{\log n}{\varepsilon^2})\).

Note that, to guarantee the correctness of the algorithm, the sampling cannot be arbitrary; it must satisfy the condition of Lemma~\ref{lem:small-supp-approx} below.

The paper gives two lemmas:

\begin{lemma}[Small Support Approximation Lemma {\cite[Lemma~1]{12WSNE} (implied by the proof of Lemma \ref{lem:lmm_thm} \cite[Theorem~1]{lmm03})}]
  \label{lem:small-supp-approx}
  Let \((\mathbf{x}, \mathbf{y})\) be a strategy profile of an \(n \times n\) bimatrix game. Then there exists an \(O(\log n / \varepsilon^2)\)-uniform strategy profile \((\mathbf{x}_s, \mathbf{y}_s)\) sampled from \((\mathbf{x}, \mathbf{y})\) such that
  \begin{align*}
    |\mathcal{R}(\mathbf{x}_s, \mathbf{y}_s) - \mathcal{R}(\mathbf{x}, \mathbf{y})| & \le \varepsilon                        \\
    |\mathcal{C}(\mathbf{x}_s, \mathbf{y}_s) - \mathcal{C}(\mathbf{x}, \mathbf{y})| & \le \varepsilon                        \\
    |\mathcal{R}(e_i, \mathbf{y}_s) - \mathcal{R}(e_i, \mathbf{y})|                 & \le \varepsilon &  & \forall i \in [n] \\
    |\mathcal{C}(\mathbf{x}_s, e_j) - \mathcal{C}(\mathbf{x}, e_j)|                 & \le \varepsilon &  & \forall j \in [n]
  \end{align*}
\end{lemma}

Here ``sampled from \((\mathbf{x}, \mathbf{y})\)'' means that pure strategies are drawn according to \((\mathbf{x}, \mathbf{y})\), so strategies outside their supports are never selected. The proof uses the sampling and concentration argument from the LMM lemma above. The key estimate for the first two conditions is the decomposition
\[
  |\mathcal{R}(\mathbf{x}_s, \mathbf{y}_s) - \mathcal{R}(\mathbf{x}, \mathbf{y})| \le |\mathcal{R}(\mathbf{x}_s, \mathbf{y}) - \mathcal{R}(\mathbf{x}, \mathbf{y})| + \max_{i} |\mathcal{R}(e_i, \mathbf{y}_s) - \mathcal{R}(e_i, \mathbf{y})|
\]
and its column-player analogue. Hoeffding's inequality and a union bound over the \(2n+2\) conditions then give \(k=O(\ln n/\varepsilon^2)\); choosing a larger hidden constant makes the failure probability arbitrarily small.

\begin{lemma}[\(3\varepsilon\)-WSNE Sampling Lemma {\cite[Lemma~2]{12WSNE}}]
  \label{lem:3_eps-wsne-sampling}
  Let \((\mathbf{x}, \mathbf{y})\) be an \(\varepsilon\)-WSNE of an \(n \times n\) bimatrix game, and let \((\mathbf{x}_s, \mathbf{y}_s)\) be an \(O(\log n / \varepsilon^2)\)-uniform strategy profile satisfying the condition of Lemma~\ref{lem:small-supp-approx}. Then \((\mathbf{x}_s, \mathbf{y}_s)\) is a \(3\varepsilon\)-WSNE of the game.
\end{lemma}

The proof is given in Appendix~\ref{app:sampled-wsne-proof}. The randomized implementation and its failure probability are specified in Section~\ref{sec:source-corrections}.

\subsubsection{Preprocessing}

\begin{enumerate}
  \item With \(O\left(\frac{n \log n}{\varepsilon^4}\right)\) queries, obtain two \(\varepsilon\)-WSNEs \((\mathbf{x}^*, \mathbf{y}^*)\) and \((\hat{\mathbf{x}}, \hat{\mathbf{y}})\);
  \item Draw \(\mathbf{x}_s^*, \mathbf{y}_s^*\) with the sample size from Section~\ref{sec:source-corrections}, query all \((e_i, e_j)\) with \(i \in \supp(\mathbf{x}_s^*)\), and construct the restricted game \((\mathcal{R}_{\mathbf{x}_s^*}, \mathcal{C}_{\mathbf{x}_s^*})\).
\end{enumerate}

Then, following Algorithm 1, we case-analyze, again taking \(\mathcal{R}(\mathbf{x}^*,\mathbf{y}^*) \ge \mathcal{C}(\hat{\mathbf{x}},\hat{\mathbf{y}})\) as an example:

\subsubsection{The \texorpdfstring{\(\mathcal{R}(\mathbf{x}^*, \mathbf{y}^*) \ge \mathcal{C}(\hat{\mathbf{x}}, \hat{\mathbf{y}})\)}{R(x*,y*) >= C(xh,yh)} branch}

\paragraph{Case 2(a): \(\mathcal{R}(\mathbf{x}^*,\mathbf{y}^*) \le 1/2\)}
Simply output \((\hat{\mathbf{x}}, \mathbf{y}^*)\).

\paragraph{Case 2(b): ``one high, one low''}
In this case, \(\mathcal{R}(\mathbf{x}^*,\mathbf{y}^*) > 1/2\), and there exists \(\mathbf{x}'\) supported on \(\supp(\mathbf{x}_s^*)\) with \(\mathcal{C}(\mathbf{x}', e_j) \le 1/2\); the output is \((\mathbf{x}', \mathbf{y}_s^*)\).

\paragraph{Case 2(c): ``both high''}
In this case, neither of the previous two conditions holds.

Consider the restricted game \((\mathcal{R}_{\mathbf{x}_s^*}, \mathcal{C}_{\mathbf{x}_s^*})\). One can show that every Nash equilibrium \((\mathbf{x}, \mathbf{y})\) of \((\mathcal{R}_{\mathbf{x}_s^*}, \mathcal{C}_{\mathbf{x}_s^*})\) satisfies:
\[
  \mathcal{R}(\mathbf{x}, \mathbf{y}) > \frac12 - 3\varepsilon,\;\; \mathcal{C}(\mathbf{x}, \mathbf{y}) > \frac12 - 3\varepsilon
\]

Correctness proofs of the three cases are given in Appendix~\ref{app:case-two-low},~\ref{app:case-two-low-high}, and \ref{app:case-two-both-high}, respectively.

Enumerating all \(\kappa(\delta)\)-uniform strategy profiles, we find a \((1/2+3\varepsilon+\delta)\)-WSNE. The time complexity analysis is the same as in Algorithm 1.
